\documentclass[12pt]{article}
\usepackage[T1]{fontenc}
\usepackage[utf8]{inputenc}
\usepackage{lmodern}
\usepackage{textcomp}
\usepackage{enumitem}
\usepackage{geometry}
\usepackage{graphicx}
\usepackage{amsmath, amssymb}
\geometry{margin=1in}
\begin{document}
\begin{center}
Chemistry - Graduate Level Assessment 15 \
Total Marks: 40 \
Answer all questions.
\end{center}

\section*{Multiple Choice (10 marks)}
\begin{enumerate}[label=\arabic*.]
\item If$_{1}$.0 g of rubidium and 1.0 g of bromine are reacted, what will be left in measurable amounts (more than 0.10 mg) in the reaction vessel?
\begin{enumerate}[label=(\alph*)]
    \item RbBr only
    \item RbBr, Rb, Br$_{2}$, and Rb$_{2}$Br
    \item RbBr and Rb$_{2}$Br only
    \item RbBr, Rb, and Br$_{2}$
    \item Rb and Br$_{2}$ only
\end{enumerate}
\item Cerium(III) sulfate, Ce$_{2}$(SO 4)2, is less soluble in hot water than it is in cold. Which of the following conclusions may be related to this?
\begin{enumerate}[label=(\alph*)]
    \item The hydration energies of cerium ions and sulfate ions are very high.
    \item The heat of solution of cerium(III) sulfate is isothermic.
    \item The solubility of cerium(III) sulfate increases with temperature.
    \item The heat of solution of cerium(III) sulfate is exothermic.
    \item The hydration energies of cerium ions and sulfate ions are very low.
\end{enumerate}
\item Calculate the ratio of the electrical and gravitational forces between a proton and an electron.
\begin{enumerate}[label=(\alph*)]
    \item 4 $10^{39}$
    \item 6 $10^{36}$
    \item 7 $10^{37}$
    \item 5 $10^{39}$
    \item 9 $10^{41}$
\end{enumerate}
\item When 10.0 g of n-propylchloride is allowed to react with excess sodium in the Wurtz reaction, how many grams of hexane would be produced assuming a 70\% yield?
\begin{enumerate}[label=(\alph*)]
    \item 8.67 g
    \item 3.82 g
    \item 5.11 g
    \item 3.00 g
    \item 4.45 g
\end{enumerate}
\item What is the molar mass of a monoprotic weak acid that requires 26.3 mL of 0.122 M KOH to neutralize 0.682 gram of the acid?
\begin{enumerate}[label=(\alph*)]
    \item 26.3 g mol-1
    \item 500 g mol-1
    \item 212 g mol-1
    \item 0.122 g mol-1
    \item 1.22 g mol-1
\end{enumerate}
\end{enumerate}

\section*{True / False (10 marks)}
\textit{Answer True or False}
\begin{enumerate}[label=\arabic*.]
\setcounter{enumi}{5}
\item True or False: The van der Waals parameter a is measured in SI base units as 7.61 \ensuremath{\times} 10^\{-2\} kg $m^5$ s^\{-2\} mol^\{-2\}.
\item True or False: One liter of paint, brushed out to a thickness of 100 microns, will cover an area of 30 square meters.
\item True or False: The mass of water vapor present in a 400 $m^3$ room at 27^\ensuremath{\circC} and 60\% relative humidity is 6.2 kg.
\item True or False: The thermal stability of group-14 hydrides decreases in the order CH 4 \textless{} SiH 4 \textless{} GeH 4 \textless{} SnH 4 \textless{} PbH 4.
\item True or False: Sulfur dioxide (SO 2) is classified as an acid anhydride due to its ability to form sulfurous acid upon hydration.
\end{enumerate}

\section*{Long Form Response (20 marks)}
\textit{Answer each question with a clear and complete explanation.}
\begin{enumerate}[label=\arabic*.]
\setcounter{enumi}{10}
\item Calculate the frequency of a light wave that corresponds to the wavelength equal to the Bohr radius (0.5292 \ensuremath{\times} 10^-8 cm). Discuss the implications of this relationship in quantum mechanics.
\item Discuss the criteria for Hermitian operators in quantum mechanics and analyze why the operators d/dx and $d^2$/$dx^2$ qualify as Hermitian, including their implications in chemical systems.
\end{enumerate}

\vfill
\noindent\textit{End of Paper}
\end{document}